\section{指标计算}
\label{sec:ch15-index-calculation}

本节对相关的\personRiemann{}--\personHilbert{}问题作指标计算。

回忆，我们固定了\termBasePath{} $\ell_0:[0,1]\to M$，并考虑带基点道路空间
$\Omega(L_0,L_1;\ell_0)$。还取 $\ell_0^*\Lambda^{\mathrm{ori}}(TM)$ 的截面
$t\mapsto\lambda_0(t)$，使
\[
 \lambda_0(0)=T_{\ell_0(0)}L_0,
 \qquad
 \lambda_0(1)=T_{\ell_0(1)}L_1.
\]
对 $p\in L_0\cap L_1$，考虑拉格朗日子空间道路空间
\[
 \mathcal P(T_pL_0,T_pL_1)
 :=\{\lambda:[0,1]\to\Lambda(T_pM)\mid
 \lambda(0)=T_pL_0,\ \lambda(1)=T_pL_1\}.
\]
利用上述道路 $\lambda_0$ 和对 $[\ell_p,w]$，按如下方式定义
$\mathcal P(T_pL_0,T_pL_1)$ 中一个元素的同伦类。因为 $[0,1]^2$ 可缩，存在同构
$w^*TM\cong[0,1]^2\times T_pM$。于是把道路 $\lambda_{w,\lambda_0}$ 定义为下列
道路的拼接：
\[
 \lambda(0,t)=\lambda_0(t),
 \qquad
 \lambda(s,i)=T_{w(s,i)}L_i.
\]

考虑半带
\[
 Z_-:=\{z\in\mathbb C\mid\operatorname{Re}z\geq0,
 0\leq\operatorname{Im}z\leq1\},
 \qquad
 Z_+:=\{z\in\mathbb C\mid\operatorname{Re}z\leq0,
 0\leq\operatorname{Im}z\leq1\},
\]
其中 $z=\tau+\sqrt{-1}\,t$。对每条
$\lambda\in\mathcal P(T_pL_0,T_pL_1)$，如下定义\personFredholm{}算子。
令 $W_\lambda^{1,p}(Z_+;T_pM)$（分别地，$W_\lambda^{1,p}(Z_-;T_pM)$）
为由 $L_{\mathrm{loc}}^{1,p}$ 映射 $\zeta_+:Z_+\to T_pM$（分别地，
$\zeta_-:Z_-\to T_pM$）组成的\personBanach{}空间，并要求
\begin{enumerate}[label=(\arabic*)]
\item $\zeta_+(\tau,i)\in T_pL_i$（分别地，
$\zeta_-(\tau,i)\in T_pL_i$），其中 $i=0,1$；
\item $\zeta_+(0,t)\in\lambda(t)$（分别地，
$\zeta_-(0,t)\in\lambda(t)$）。
\end{enumerate}
\footnote{译者注：原书第 (2) 项把负半带条件印作
$\zeta_-(z_-)\in\lambda(t)\ (i=0,1)$；其中 $z_-$ 未定义，尾部 $i=0,1$
也与变量 $t$ 不相干。由正半带条件及半带的竖直边界，译文补正为
$\zeta_-(0,t)\in\lambda(t)$。}
令 $L^p(Z_\pm;\Lambda^{0,1}(Z_\pm)\otimes T_pM)$ 为丛
$\Lambda^{0,1}(Z_\pm)\otimes T_pM$ 的 $L_{\mathrm{loc}}^p$ 截面 $\zeta_\pm$
组成的\personBanach{}空间。\personDolbeault{}算子诱导有界线性映射
\[
 \overline\partial_{\lambda,Z_\pm}:
 W_\lambda^{1,p}(Z_\pm;T_pM)
 \longrightarrow L^p(Z_\pm;\Lambda^{0,1}(Z_\pm)\otimes T_pM),
\]
它是\personFredholm{}算子。

\begin{proposition}
\label{prop:semi-strip-index-maslov-morse}
设 $\mu([\ell_p,w];(\ell_0,\lambda_0))$ 是定义~\ref{def:maslov-morse-index}
给出的\termMaslovMorseIndex{}。则
\begin{equation}
 \operatorname{Index}\overline\partial_{(\lambda_w,Z_-)}
 =\mu([\ell_p,w];(\ell_0,\lambda_0)).
 \label{eq:semi-strip-index-maslov-morse}
\end{equation}
\end{proposition}

\begin{proof}
令 $\lambda_w$ 为第~\ref{sec:maslov-morse-index}节定义的环路，并考虑\termBundlePair{}
$(w^*TM,\lambda_w)$。对 $(s,0),(s,1)\in\partial[0,1]^2$，记
\[
 \lambda_w(s,i)=T_{w(s,i)}L_i,
 \qquad i=0,1.
\]
现在把这些拉格朗日道路补成一个丛对。为此，需要取
$\Lambda(T_pM)$ 中的一条拉格朗日道路 $\alpha_p$，使
\begin{equation}
 \alpha_p(0)=T_pL_0,
 \qquad
 \alpha_p(1)=T_pL_1.
 \label{eq:reference-lagrangian-path-alpha-p}
\end{equation}
选取 $T_pM$ 中一条有向拉格朗日子空间道路 $\alpha_p=\alpha_p(t)$，使它延拓
式~\eqref{eq:reference-lagrangian-path-alpha-p}中两个拉格朗日子空间上的给定定向。
通过对 $\alpha_p$ 施加一个条件，固定其同伦类。

注意，在平凡化 $w^*TM\cong[0,1]^2\times T_pM$ 下，拉格朗日子丛
$\lambda_w$ 定义一条拉格朗日环路，仍记为 $\lambda_w$。该环路的\personMaslov{}指标
在辛变换和同伦下不变。考虑在边界
\[
 \{(1,t)\mid t\in[0,1]\}\subset\partial[0,1]^2
\]
上固定的环路同伦。随后借助拉长函数
$\rho:[0,\infty)\to[0,1]$，把
$[0,1)\times[0,1]\cong\mathbb R_+\times[0,1]$ 拉长。在此拉长下，算子
$\overline\partial_{(\lambda_w;Z_-)}$ 在 $\tau\to\infty$ 时成为柱面型；由\personFredholm{}指标的\termHomotopyInvariance{}，其指标在固定渐近条件的系数和边界条件连续
形变下不变。此类形变论证的细节见附录 C。

因而只需考虑如下特殊情形：
\[
 T_pM=\mathbb C^n,
 \qquad T_pL_0=\mathbb R^n,
 \qquad T_pL_1=\sqrt{-1}\,\mathbb R^n,
\]
以及
\[
 \alpha_p(t)=e^{\pi\sqrt{-1}t/2}\mathbb R^n,
\]
连同 $\partial[0,1]^2\setminus\{(1,t)\}\to\mathbb C^n$ 的道路
\[
\begin{aligned}
 \lambda_w(s,0)&\equiv\mathbb R^n,\\
 \lambda_w(s,1)&\equiv\sqrt{-1}\,\mathbb R^n,\\
 \lambda_w(0,t)&=
 e^{-(l_1+1/2)\pi\sqrt{-1}t}\mathbb R
 \oplus\cdots\oplus
 e^{-(l_n+1/2)\pi\sqrt{-1}t}\mathbb R,
\end{aligned}
\]
其中 $l_j\in\mathbb Z$。于是拉格朗日环路为 $\lambda_w=\lambda_w\cup\alpha_p$。

指标问题分裂为如下类型的一维问题：
\[
 \begin{cases}
 \overline\partial\zeta=0,\\
 \zeta(\tau,0)\in\mathbb R,
 \quad\zeta(\tau,1)\in\sqrt{-1}\,\mathbb R,\\
 \zeta(0,t)\in e^{-(l+1/2)\pi\sqrt{-1}t}\mathbb R.
 \end{cases}
\]
上面在 $\partial[0,1]^2$ 上定义的相应拉格朗日环路变为
\[
\begin{array}{ll}
 (s,0)\longmapsto\mathbb R,
 & (0,t)\longmapsto e^{-(l+1/2)\pi\sqrt{-1}t}\mathbb R,\\
 (s,1)\longmapsto\sqrt{-1}\,\mathbb R,
 & (1,t)\longmapsto e^{\pi\sqrt{-1}t/2}\mathbb R.
\end{array}
\]
这条沿 $\partial[0,1]^2$ 正向行进的环路，同伦于如下拼接的拉格朗日道路：
先取
\[
 t\in[0,1]\longmapsto e^{(l+1/2)\pi\sqrt{-1}(1-t)}\mathbb R,
\]
再取
\[
 t\in[0,1]\longmapsto e^{\pi\sqrt{-1}t/2}\mathbb R.
\]
该拼接道路定义一条\personMaslov{}指标为 $l+1$ 的环路。求和得到
\[
 \mu(\lambda_w)=\sum_{j=1}^n(l_j+1).
\]
另一方面，对应的\personFredholm{}指标可由单复变量\personFourier{}分析显式求出。细节仍见
附录 C 的命题 C.4.1（另见 (Oh99) 的 (A.12)）。这证明
\begin{equation}
 \operatorname{Index}\overline\partial_{(\lambda_w,Z_-)}
 =\sum_{j=1}^n(l_j+1).
 \label{eq:semi-strip-index-explicit}
\end{equation}
证明完成。
\end{proof}

有了这一指标公式，现在可以计算\termFloerModuliSpace{}的\termVirtualDimension{}。
下面是本节的主要定理。

\begin{theorem}
\label{thm:floer-moduli-dimension-maslov-difference}
设 $B_{ww'}\in\pi_2(p,q)$ 是满足 $w\#B_{ww'}\sim w'$ 的同伦类。则
\begin{equation}
 \dim\widetilde{\mathcal M}
 (L_0,L_1;[\ell_p,w],[\ell_q,w'])
 =\mu([\ell_p,w])-\mu([\ell_q,w'])
 =\mu(p,q;B_{ww'}).
 \label{eq:floer-moduli-dimension-index}
\end{equation}
\end{theorem}

\begin{proof}
证明基于\personFredholm{}指标的粘合公式
\begin{equation}
 \operatorname{Index}\overline\partial_{(\lambda_w,Z_-)}
 +\operatorname{Index}D_u\overline\partial
 =\operatorname{Index}\overline\partial_{(\lambda_{w'},Z_-)}.
 \label{eq:fredholm-index-gluing-ch15}
\end{equation}
其证明参见 (APS75) 定理 3.10 的证明。把它改写为
\[
 \operatorname{Index}D_u\overline\partial
 =\operatorname{Index}\overline\partial_{(\lambda_{w'},Z_-)}
 -\operatorname{Index}\overline\partial_{(\lambda_w,Z_-)}.
\]
再应用命题~\ref{prop:semi-strip-index-maslov-morse}，便得到
式~\eqref{eq:floer-moduli-dimension-index}，证明完成。
\end{proof}
